\section{New rational constraints from subset scaling}
\label{sec:subset-scaling-new}

The following arguments supply necessary conditions for global maximizers
without assuming any polygonal or reflection symmetry. They improve the
separation and radius constants in the existing compact formulation, but
do not by themselves classify the maximizing shapes for arbitrary $n$.
All proofs below are self-contained; no priority assertion is intended.
Put $N=n(n-1)/2$ and $\lambda=(n-1)/2$.

\begin{lemma}[A boundary maximum derivative]\label{lem:boundary-derivative-new}
Let $q$ be a nonconstant polynomial with
$|q(1)|=\max_{|t|=1}|q(t)|$. Then $q(1)\ne0$ and
\[
                 \Re\frac{q'(1)}{q(1)}>0.
\]
\end{lemma}
\begin{proof}
Set $f=q/q(1)$ and $b=f(0)$. The maximum-modulus principle gives
$|f|\le1$ in the unit disk, and $|b|<1$ since $f$ is nonconstant.
Schwarz--Pick, applied to $f$ at zero, gives for $0<r<1$
\[
 \left|\frac{f(r)-b}{1-\overline b f(r)}\right|\le r.
\]
Using the elementary identity for this disk automorphism yields
\[
 1-|f(r)|^2\ge
 (1-r^2)\frac{|1-\overline b f(r)|^2}{1-|b|^2}.
\]
Divide by $2(1-r)$ and let $r\uparrow1$. Since $f(1)=1$ and $f$
is analytic across the boundary, this gives
\[
 \Re f'(1)\ge\frac{|1-b|^2}{1-|b|^2}>0.
\]
\end{proof}

\begin{theorem}[Every-subset internal moment]\label{thm:subset-spread-new}
Let $z_1,\ldots,z_n$ be a global maximizer with
$\sum_i|z_i|^2=n$. For a subset $S$ of size $m$, put
\[
 a_S=\frac1m\sum_{i\in S}z_i,\qquad
 Q_S=\sum_{i\in S}|z_i-a_S|^2.
\]
For every proper subset with $2\le m<n$,
\begin{equation}\label{eq:subset-spread-new}
             Q_S>\frac{m(m-1)}{n-1}.
\end{equation}
For the full set $Q_S=n$, and for a singleton $Q_S=0$.
\end{theorem}
\begin{proof}
A maximizer is centered and collision-free. Homogeneity shows that it
also globally maximizes
\[
 \Phi(z)=\Delta_n(z)\exp\left(-\lambda\sum_i|z_i|^2\right)
\]
over all of $\C^n$: for a fixed nonzero shape of moment $Q$, the best
scale satisfies $s^2Q=N/\lambda=n$.
Fix a proper subset $S$, all outside points, and its centroid $a_S$.
Write $y_i=z_i-a_S$, so $\sum_{i\in S}y_i=0$, and replace the subset
by $a_S+t y_i$. Up to a positive factor independent of $t$, the
conditional objective is
\[
 |P(t)|^2e^{-\alpha|t|^2},\qquad
 P(t)=t^k q(t),\qquad k=\frac{m(m-1)}2,\qquad
 \alpha=\lambda Q_S,
\]
where
\[
 q(t)=\prod_{i\in S,\ j\notin S}(a_S+t y_i-z_j).
\]
The original $t=1$ globally maximizes this expression. On the unit
circle the Gaussian factor and $|t|^{2k}$ are constant, so
$|q(1)|=\max_{|t|=1}|q(t)|$. Stationarity at $t=1$ gives
\[
 \frac{P'(1)}{P(1)}=\alpha,
 \qquad \frac{q'(1)}{q(1)}=\alpha-k.
\]
The polynomial $q$ is nonconstant: some $y_i\ne0$ because $m\ge2$
and the points are distinct; its outside factors have positive degree.
If a $y_i=0$, its outside factors are nonzero constants by
collision-freeness. Thus the degree cannot cancel in the product.
Lemma~\ref{lem:boundary-derivative-new} yields $\alpha>k$, proving
\eqref{eq:subset-spread-new}. The full-set and singleton statements
follow from centering and the prescribed moment.
\end{proof}

\begin{remark}
For closed rational optimization constraints one may relax the strict
inequality to $Q_S\ge m(m-1)/(n-1)$. This weaker version also follows
using only the usual maximum-modulus principle: if
$M_q(r)=\max_{|t|=r}|q(t)|$, then $M_q$ is nondecreasing, while
$r^{2k}e^{-\alpha r^2}$ is strictly increasing for
$r<\sqrt{k/\alpha}$. Consequently a global radial maximizer at $r=1$
requires $\alpha\ge k$.
\end{remark}

\begin{corollary}[Pair and triple separation]\label{cor:pair-triple-new}
For $n\ge3$ every pair in a global maximizer satisfies
\begin{equation}\label{eq:pair-new-rational}
                    |z_i-z_j|^2>\frac4{n-1}.
\end{equation}
For $n\ge4$ every triple $S$ satisfies
\begin{equation}\label{eq:triple-new-rational}
 Q_S>\frac6{n-1},\qquad
 \sum_{\{i,j\}\subset S}|z_i-z_j|^2>\frac{18}{n-1}.
\end{equation}
At force-one scale $w_i=\sqrt\lambda\,z_i$, these become
$|w_i-w_j|^2>2$, and respectively $Q_S(w)>3$ and a triple
squared-distance sum greater than $9$.
\end{corollary}
\begin{proof}
For any $m$ points,
$mQ_S=\sum_{\{i,j\}\subset S}|z_i-z_j|^2$.
Apply Theorem~\ref{thm:subset-spread-new} for $m=2,3$ and rescale.
For $n=2$ the pair squared distance is $4$ in the original scale and
$2$ at force-one scale. Since $e>2$, \eqref{eq:pair-new-rational}
is strictly stronger than $8/[e(n-1)]$ whenever $n\ge3$.
\end{proof}

\begin{corollary}[Radius and every-subset centroid]\label{cor:subset-centroid-new}
For $n\ge3$ and every nonempty proper subset $S$ of size $m$,
\begin{equation}\label{eq:subset-centroid-new}
 \left|\sum_{i\in S}z_i\right|^2
      <\frac{2m^2(n-m)^2}{n(n-1)}.
\end{equation}
In particular,
\begin{equation}\label{eq:radius-new-rational}
                    |z_i|^2<\frac{2(n-1)}n.
\end{equation}
At force-one scale the bounds are
\[
 \left|\sum_{i\in S}w_i\right|^2
       <\frac{m^2(n-m)^2}{n},\qquad
 |w_i|^2<\rho_n:=\frac{(n-1)^2}{n}.
\]
All displayed strict inequalities may be relaxed to nonstrict rational
constraints retaining every global maximizer. For $n=2$ the radius
inequality is equality.
\end{corollary}
\begin{proof}
Let $T$ be the complement. Centering gives
\[
 n=Q_S+Q_T+\frac{n}{m(n-m)}
                    \left|\sum_{i\in S}z_i\right|^2.
\]
Theorem~\ref{thm:subset-spread-new} gives
\[
 Q_S+Q_T>
 \frac{m(m-1)+(n-m)(n-m-1)}{n-1}
 =n-\frac{2m(n-m)}{n-1}.
\]
Singleton blocks have internal moment zero; at least one block has
size at least two when $n\ge3$, so the sum is strict. Rearrangement
proves \eqref{eq:subset-centroid-new}, and $m=1$ gives
\eqref{eq:radius-new-rational}.
\end{proof}

\begin{proposition}[Directional majorization]\label{prop:projection-new}
Fix an angle $\theta$, and sort
$x_i=\Re(e^{-i\theta}z_i)$ in decreasing order. Put
$c=\sqrt{2/[n(n-1)]}$ and $b_i=c(n+1-2i)$. Then $(x_i)$ is
majorized by $(b_i)$: for $1\le m<n$,
\[
 \sum_{i=1}^m x_i\le c m(n-m)=\sum_{i=1}^m b_i,
 \qquad \sum_i x_i=\sum_i b_i=0.
\]
For $n\ge3$ every proper partial-sum inequality is strict. Consequently,
for every convex real function $\varphi$,
\[
                   \sum_i\varphi(x_i)\le\sum_i\varphi(b_i).
\]
In particular, for $n\ge3$,
\begin{equation}\label{eq:quadrupole-new}
 \sum_i x_i^2<\frac{2(n+1)}3,\qquad
 \left|\sum_i z_i^2\right|<\frac{n+4}3.
\end{equation}
For real $t\ne0$,
\[
 \sum_i e^{t x_i}\le\frac{\sinh(nct)}{\sinh(ct)}.
\]
The ratio has continuous value $n$ at $t=0$.
\end{proposition}
\begin{proof}
Apply Corollary~\ref{cor:subset-centroid-new} to the subset giving the
largest $m$ projections; the projection of its sum is bounded by the
modulus of that sum. The arithmetic progression $(b_i)$ has the stated
partial sums, proving majorization. The convex-function inequality is
the elementary majorization inequality. For the square it is strict
because strict proper partial sums exclude a permutation of $(b_i)$.
Now
\[
 \sum_i b_i^2=\frac{2(n+1)}3,\qquad
 \max_\theta\sum_i\Re(e^{-i\theta}z_i)^2
       =\frac12\left(n+\left|\sum_i z_i^2\right|\right),
\]
which proves \eqref{eq:quadrupole-new}. The exponential expression is
just the finite geometric sum over the arithmetic progression.
\end{proof}

\begin{remark}[Population cut]
If a disk of radius $r$ contains $m$ points, then its internal moment
is at most $mr^2$. For a proper subset with $m\ge2$, the new theorem
therefore gives $m-1<(n-1)r^2$. This is a useful finite population cut,
but its coefficient is still twice the limiting bulk count in the
uniform disk of radius $\sqrt2$. No density classification or exact
all-$n$ formula follows from this cut alone.
\end{remark}

\subsection{Rational trace bounds using coefficient Euclidean norms}
Write the force-one polynomial as
$P(t)=\prod_i(t-w_i)=t^n+c_1t^{n-1}+\cdots+c_n$, and let
$V_{ji}=w_i^j$ for $0\le j<n$. On physical roots the Hermitian trace
matrix is $H=VV^*$. Put $d=n-1$, $\rho_n=(n-1)^2/n$, and define
\[
 T_{n,k}=\binom dk^2\rho_n^k\quad(0\le k\le d),\qquad
 A_n=\sum_{k=0}^dT_{n,k},\qquad T_n=\max_{0\le k\le d}T_{n,k}.
\]
All these constants are positive rationals.

\begin{proposition}[Improved compact coefficient constants]\label{prop:trace-new}
Every global maximizer at force-one scale satisfies
\begin{equation}\label{eq:trace-new}
 H\succeq\mu_n I,\qquad
 \mu_n=\frac{2^d}{nA_n},\qquad
 |c_k|^2\le\binom nk^2\rho_n^k.
\end{equation}
Let $R(t)=\sum_{\ell=0}^{n-1}r_\ell t^\ell$ be the unique polynomial
of degree below $n$ satisfying $R(w_i)=\overline w_i$. Then
\begin{equation}\label{eq:conjugation-coeff-new}
 |r_\ell|^2\le B_n:=\frac{NnT_n}{2^d}.
\end{equation}
At $n=6$ the constants are
\[
 \rho_6=\frac{25}6,\qquad A_6=\frac{138927151}{7776},\qquad
 \mu_6=\frac{41472}{138927151}>\frac1{3350},\qquad
 B_6=\frac{48828125}{2304}<21193.
\]
\end{proposition}
\begin{proof}
The radius bound is $|w_i|^2\le\rho_n$, and the pair bound is
$|w_i-w_j|^2\ge2$. Hence
\[
 |P'(w_i)|^2=\prod_{j\ne i}|w_i-w_j|^2\ge2^d.
\]
For the Lagrange polynomial
\[
 L_i(t)=\frac{\prod_{j\ne i}(t-w_j)}{P'(w_i)},
\]
the coefficient of $t^{d-k}$ has squared modulus at most
$T_{n,k}/2^d$. Thus its squared coefficient Euclidean norm is at most
$A_n/2^d$. These coefficient vectors form the inverse of $V^T$;
therefore
\[
 \|V^{-1}\|_F^2\le\frac{nA_n}{2^d},\qquad
 \lambda_{\min}(H)\ge\|V^{-1}\|_F^{-2}\ge\mu_n.
\]
The elementary symmetric functions give the bound on $c_k$.
Finally, interpolation gives $R=\sum_i\overline w_i L_i$.
For each coefficient, Cauchy--Schwarz and $\sum_i|w_i|^2=N$ give
\[
 |r_\ell|^2\le
 \left(\sum_i|w_i|^2\right)
 \left(\sum_i|[t^\ell]L_i|^2\right)
 \le\frac{NnT_n}{2^d}.
\]
The six-point values follow by rational arithmetic.
\end{proof}

The disk localizing matrix can now use
\[
 J_{jk}=\rho_n H_{jk}-H_{j+1,k+1}
       =\bigl[V\operatorname{diag}(\rho_n-|w_i|^2)V^*\bigr]_{jk}.
\]
For the separation-product cut, the quotient-algebra polynomial
$D(t)=P'(t)\overline{P'}(R(t))$ satisfies
\[
 K_{jk}=\operatorname{Tr}\bigl(t^jR(t)^k[D(t)-2^d]\bigr)
       =\bigl[V\operatorname{diag}(|P'(w_i)|^2-2^d)V^*\bigr]_{jk}.
\]
Hence both $J\succeq0$ and $K\succeq0$ are closed rational necessary
conditions for every global maximizer. As with the existing model,
$K$ bounds products of separations and does not independently impose
every pair inequality equivalently. These strengthened cuts preserve
the exact maximum when added to a model whose physical points are
actual moment-$n$ configurations; they do not turn necessary conditions
into a classification theorem.
